\documentclass[11pt]{article}
\usepackage[a4paper,margin=25mm]{geometry}
\usepackage{amsmath,amssymb,amsthm,array,booktabs,longtable}
\usepackage[T1]{fontenc}
\usepackage{lmodern}
\usepackage[hidelinks]{hyperref}
\newtheorem{theorem}{Theorem}
\newtheorem{corollary}[theorem]{Corollary}
\newcommand{\floor}[1]{\left\lfloor #1\right\rfloor}
\newcommand{\ceil}[1]{\left\lceil #1\right\rceil}
\title{Correction of Kim--Stolarsky's eventual\newline three-prefix conjecture}
\author{Haobo Ma\\\small ChronoAI Pte Ltd, Singapore
\and Wenlin Zhang\\\small National University of Singapore\\\small ChronoAI Pte Ltd, Singapore}
\date{November 17, 2026}
\begin{document}
\maketitle
\begin{abstract}
Kim and Stolarsky's eventual three-prefix conjecture in \emph{Integers} 26
(2026), A98, Remark~3, p.~6, fails in both its squared and third terms whenever
$n\equiv2\pmod3$, by exact continued-fraction divisions and polynomial sign bounds.
We give corrected prefixes for all three root branches, prove the proposed
ordering and all three limiting digit ratios, and determine the full rational
expansions for the squared odd and even terms. A specialization of the authors'
own identity also answers their question about the separation of double zeros:
there is no finite bound, even for integer-coefficient factors.
\end{abstract}

\section{The source assertion and its correction}
Put $T_j=j(j+1)/2$. For integers $s\ge2$, Kim and Stolarsky's
Theorem~1~\cite{KS} defines
\begin{align*}
 L_{s,o}(x)&=\sum_{j=1}^{s}(2j-1)^3+x^3+
       \sum_{j=2}^{s}\bigl(2j-2+(j-1)x\bigr)^3,\\
 R_{s,o}(x)&=\bigl(T_{2s-1}+(1+T_{s-1})x\bigr)^2,\\
 P_{m,o}(x)&=\sum_{j=0}^{m}(4j+1+jx)
                  \bigl((j+1)(j-2)+j^2x\bigr),\\
 Q_m(x)&=xP_{m,o}(x)-x^2+x^3.
\end{align*}
Their identity is $L_{m+1,o}-R_{m+1,o}=Q_m$.
For $m=m_n=10^{4n+2}$, $n=1,2,\ldots$, their Remark~3 conjectures,
``For large $n$'', the following continued-fraction prefixes for the
negatives of roots tending to $-4$:
\begin{align}
 y_L(m)&=\left[4;\frac{3m}{4},3,4,1,\frac{9m-92}{404},\ldots\right],\label{eq:sourceL}\\
 y_R(m)&=\left[4;\frac m2+2,\floor{m/7}-1,\ldots\right],\label{eq:sourceR}\\
 y_Q(m)&=\left[4;\frac{3m}{4},3,\ceil{(6m+100)/35},\ldots\right].\label{eq:sourceQ}
\end{align}
Here $-y_L,-y_R,-y_Q$ are roots of $L_{m+1,o},R_{m+1,o},Q_m$,
respectively. The $m+1$ shifts are essential. The same conjecture occurs
in arXiv:2511.15133v1, Remark~4, p.~6. A prefix means the successive
floor-and-inverse digits of an ordinary simple continued fraction;
for rational numbers we use the expansion with final digit greater than one
(notation as in~\cite{DLMF}).

\begin{theorem}\label{thm:all}
For $m=m_n$, $n\ge1$, the three polynomials have unique roots below $-4$.
Their negatives satisfy~\eqref{eq:sourceL} and the corrected prefixes
\begin{equation}\label{eq:correct}
 y_R(m)=\left[4;\frac m2+2,\floor{(m-3)/7},\ldots\right],\qquad
 y_Q(m)=\left[4;\frac{3m}{4},3,\floor{(6m+128)/35},\ldots\right].
\end{equation}
Both original prefixes~\eqref{eq:sourceR} and~\eqref{eq:sourceQ} fail
exactly for $n\equiv2\pmod3$. Moreover,
\begin{equation}\label{eq:order}
 y_R(m)>y_L(m)>y_Q(m)>4,
\end{equation}
and the roots tend to $-4$. The ratios of the first partial quotient after
the integer part to the next large partial quotient tend, respectively, to
$101/3$, $7/2$, and $35/8$ for $L,R,Q$.
\end{theorem}

The correction is necessary at arbitrarily large indices, so no common
eventual threshold makes the three original prefixes true. The squared
term alone already gives this obstruction; the third term fails on the
same progression. The first term, ordering and all three proposed ratio
limits do hold.

\section{The squared odd term: complete expansions}
\begin{theorem}\label{thm:R}
For every even $m\ge12$,
\[
 y_R(m)=\frac{4m^2+6m+2}{m^2+m+2},\qquad
 y_R(m)=\left[4;\frac m2+2,\floor{(m-3)/7},\ldots\right].
\]
Termination occurs at the third digit if $(m-3)/7$ is integral.
For $m=m_n$, $n\ge0$, put $A=m/2+2$ and $b=\floor{m/7}$.
The full canonical expansions are
\begin{equation}\label{eq:Rfull}
 y_R(m)=\begin{cases}
 [4;A,b-1,1,6],&n\equiv0\pmod3,\\
 [4;A,b-1,1,2,2],&n\equiv1\pmod3,\\
 [4;A,b,7],&n\equiv2\pmod3.
 \end{cases}
\end{equation}
\end{theorem}
\begin{proof}
The square $R_{m+1,o}(x)$ vanishes precisely when
$T_{2m+1}+(1+T_m)x=0$, giving the displayed value of $y_R$. Write $D=m^2+m+2$.
Two divisions give
\begin{equation}\label{eq:Rdivide}
 y_R=4+\frac{2m-6}{D},\qquad
 \frac1{y_R-4}=\frac m2+2+\frac7{m-3}.
\end{equation}
For even $m\ge12$, both displayed remainders lie strictly between $0$ and
$1$, so the next complete quotient is $(m-3)/7$.
For $m=m_n$, its residue modulo $7$ is $2,1,4$ as $n\equiv0,1,2\pmod3$,
because $m_n\equiv2\cdot4^n\pmod7$. Thus that complete quotient is
$b-1+6/7$, $b-1+5/7$, or $b+1/7$, respectively.
The terminal divisions $7/6=[1;6]$, $7/5=[1;2,2]$ and $7=[7]$
prove~\eqref{eq:Rfull}. All final digits exceed one.
For any proposed threshold $N$, the index $3N+2\ge N$ has actual third
digit $b$, instead of $b-1$. This negates the eventual conjunction.
\end{proof}

\section{The cubic branches: exact prefix bounds}
We prove the remaining parts of Theorem~\ref{thm:all} by rational sign
bounds, not numerical root approximations. Set
\[
 U=T_m^2,\quad B=1+T_m,\quad C=4m^3+9m^2+6m-63,
\]
and define
\begin{align}
 F_L(z)&=(U+1)z^3+(6U+12)z^2+(12U+48)z-C,\label{eq:FL}\\
 F_Q(z)&=F_L(z)+(m-3-Bz)^2.\label{eq:FQ}
\end{align}
Power-sum formulas give $F_L(z)=-L_{m+1,o}(-4-z)$; the source identity gives
$F_Q(z)=-Q_m(-4-z)$.
For $m\ge12$, $F_L$ is strictly increasing on $z\ge0$ and $F_L(0)<0$.
For $F_Q$, the coefficients of $z^3,z^2,z$ are positive; the latter two are
\[
 7U+m^2+m+13,\qquad
 3m^4+5m^3+5m^2+m+54,
\]
The first of these is $6U+12+B^2$.
Its constant coefficient is $-4(m-2)(m^2+4m+9)<0$.
Consequently each has exactly one positive zero, denoted $z_L,z_Q$.

For $t>1$, the two rational functions
\begin{equation}\label{eq:Z}
 Z_L(t)=\frac{4(16t+13)}{48mt+39m+20t+16},\qquad
 Z_Q(t)=\frac{4(3t+1)}{9mt+3m+4t}
\end{equation}
satisfy
\[
 4+Z_L(t)=[4;3m/4,3,4,1,t],\qquad
 4+Z_Q(t)=[4;3m/4,3,t].
\]
Both are strictly decreasing in $t$, with derivatives $-16$ divided by
the respective squared denominators. Throughout the prefix argument,
$m\ge M:=10^6$.

Put $d=(9m-92)/404$. Substitution in~\eqref{eq:FL} gives
\begin{align}
 Z_L(d)&=\frac{36m+945}{27m^2+720m+289},&
 4(27m^2+720m+289)^3F_L(Z_L(d))&=P_0(m),\label{eq:LP0}\\
 Z_L(d+1)&=\frac{36m+2561}{27m^2+1932m+794},&
 4(27m^2+1932m+794)^3F_L(Z_L(d+1))&=-P_1(m).\label{eq:LP1}
\end{align}
The coefficients of $P_0,P_1$ are given in Appendix~\ref{app:signs}.
Every coefficient of $P_0$ is positive. Also
$P_1(m)/m^6>133974-336>0$ for $m\ge M$, by the displayed coefficient
bound there. Hence the unique zero satisfies
$z_L=Z_L(t_L)$ for $d<t_L<d+1$.
For $m=m_n$, $3m/4$ and $d$ are positive integers:
$10^4\equiv1\pmod{101}$ makes $9m-92$ divisible by $404$.
Thus~\eqref{eq:sourceL} holds, including its last digit.

For $Q$, direct substitution gives
\begin{equation}\label{eq:HQ}
 F_Q(Z_Q(t))=\frac{-8(9mt+3m+7t+1)H_m(t)}{(9mt+3m+4t)^3},
\end{equation}
where
\begin{align*}
 H_m(t)={}&(35m^3-719m^2-1134m-432)t^2\\
 &+(-6m^4+6m^3-512m^2-548m-120)t\\
 &-2m^4-m^3-89m^2-54m-8.
\end{align*}
Let $t_-=(6m+128)/35$ and $t_+=(6m+129)/35$. Exact substitution yields
\begin{align}
 1225H_m(t_-)&=-194m^4-653633m^3-16055337m^2
                 -21789398m-7625288<0,\label{eq:Hminus}\\
 1225H_m(t_+)&=16(m^4-40816m^3-1016978m^2
                 -1381575m-483782)>0.\label{eq:Hplus}
\end{align}
For the last sign divide the parenthesis by $m^4$:
the subtracted terms total less than $0.041$ for $m\ge M$.
Equations~\eqref{eq:HQ}--\eqref{eq:Hplus} and monotonicity of $F_Q,Z_Q$
place $z_Q=Z_Q(t_Q)$ at $t_-<t_Q<t_+$.
This interval lies within a single interval between consecutive integers,
with an integer possibly at an endpoint, because its endpoints have
consecutive integer numerators and denominator $35$.
Therefore $\floor{t_Q}=\floor{(6m+128)/35}$, proving~\eqref{eq:correct}.
Finally, $m_n\equiv30,15,25\pmod{35}$ for $n\equiv0,1,2\pmod3$;
comparison of the floor and ceiling shows that the corrected digit equals
the source digit in the first two classes and is one smaller in the third.

\begin{corollary}\label{cor:order}
The ordering~\eqref{eq:order} holds for every integer $m\ge12$.
The three roots approach $-4$, and the three ratio limits in
Theorem~\ref{thm:all} hold.
\end{corollary}
\begin{proof}
Write $z_R=y_R-4=(m-3)/B$. Substitution gives
\[
 F_L(z_R)=\frac{m(m-1)(m+1)^2(2m+1)
 (m^4-6m^3-39m^2-48m-12)}{(m^2+m+2)^3}>0.
\]
For the last factor put $m=12+v$: it becomes
$v^4+42v^3+609v^2+3336v+4164$, positive for $v\ge0$.
Thus $z_R>z_L$. At $z_L$,
$F_Q(z_L)=(m-3-Bz_L)^2>0$, since $z_L\ne z_R$; hence $z_L>z_Q>0$.
Moreover $z_L<C/(12U+48)\to0$, and $z_R\to0$ by~\eqref{eq:Rdivide}.
The relevant digits, divided by $m$, tend to
$(3/4,9/404)$ for $L$, $(1/2,1/7)$ for $R$, and $(3/4,6/35)$ for $Q$.
Taking their ratios proves $101/3$, $7/2$, and $35/8$.
\end{proof}

\section{Two further consequences}
\begin{corollary}\label{cor:even}
The squared even term $R_{m,e}(x)=(T_{2m}+(1+T_m)x)^2$ has unique real root
$-y_E(m)$, where
\[
 y_E(m)=\frac{4m^2+2m}{m^2+m+2}.
\]
For every even $m\ge40$,
\[
 y_E(m)=\left[3;1,\frac m2-3,1,1,\floor{(m-10)/28},\ldots\right].
\]
For $m=m_n$, $n\ge0$, put $c=\floor{m/28}$ and $E=m/2-3$.
The full canonical expansions are
\begin{equation}\label{eq:Efull}
 y_E(m)=\begin{cases}
 [3;1,E,1,1,c,4,1,2],&n\equiv0\pmod3,\\
 [3;1,E,1,1,c-1,1,13],&n\equiv1\pmod3,\\
 [3;1,E,1,1,c-1,1,3,1,2],&n\equiv2\pmod3.
 \end{cases}
\end{equation}
\end{corollary}
\begin{proof}
Set $e=(2m+8)/(m^2+m+2)$, so $y_E=4-e$. The successive complete
quotients after digits $3,1$ are
\[
 \frac{1-e}{e}=\frac m2-\frac52+\frac7{m+4},\quad
 \frac{2(m+4)}{m+18}=2-\frac{28}{m+18},\quad
 1+\frac{28}{m-10},\quad \frac{m-10}{28}.
\]
The hypotheses give the stated intervening digits. For $m=m_n$, the
residues modulo $28$ are $16,8,4$, respectively. The last quotient is
$c+3/14$, $c-1+13/14$, or $c-1+11/14$.
The expansions $14/3=[4;1,2]$, $14/13=[1;13]$ and
$14/11=[1;3,1,2]$ give~\eqref{eq:Efull}.
These are explicit even-case consequences; the source states no particular
even prefix to be refuted.
\end{proof}

\begin{corollary}\label{cor:distance}
There is no finite upper bound on the separation of two double zeros of
polynomials of the form
$\sum_{j=1}^m(a+bj+cj^2)(d+ej+fj^2)$.
This remains true for integer coefficients and nonnegative integer zeros.
\end{corollary}
\begin{proof}
On p.~4 of~\cite{KS}, the authors ask how far apart two multiple zeros
can be, and display
\[
 \sum_{j=1}^m(1-bj)(1-2j+bj^2)=-\tfrac14m^2(b-2+bm)^2.
\]
In this identity put $b=2/(k+1)$ and multiply by $(k+1)^2$, for any
integer $k\ge1$. It gives
\[
 \sum_{j=1}^m(k+1-2j)\bigl(k+1-2(k+1)j+2j^2\bigr)
       =-m^2(m-k)^2.
\]
The polynomial on the right has double zeros $0,k$, separated by $k$.
Thus the answer follows from the authors' own identity, not a new
power-sum identity.
\end{proof}

\section{Coverage and verification scope}
\begin{center}
\small
\begin{tabular}{@{}p{0.43\linewidth}p{0.53\linewidth}@{}}
\toprule
Source item & Coverage \\
\midrule
Remark 3: simultaneous eventual assertion & Refuted and corrected;
Theorem~\ref{thm:all}. \\
Remark 3: individual odd prefixes & $L$: proved. $R,Q$: refuted and
corrected; Theorems~\ref{thm:all}--\ref{thm:R}. \\
Remark 3: ordering; three ratios & Proved; Corollary~\ref{cor:order}. \\
Remark 3: unspecified even phenomena & Explicit squared-term expansions,
Corollary~\ref{cor:even}; even cubic branches not covered. \\
Introduction, p. 4: distance between multiple zeros & Answered: no finite
bound, by the authors' displayed identity; Corollary~\ref{cor:distance}. \\
Theorem 1; Remarks 1--2,4--7; Sections 4--5 & Prior identities and
separate constructions; not re-established. \\
Introduction, p. 3: full discriminant classification; conjectures in [5]
of~\cite{KS} & Not covered; the double-zero family is not a classification. \\
\bottomrule
\end{tabular}
\end{center}
The pinned
\href{https://github.com/the-omega-institute/trureturing/blob/53c631a19a8292b44b6ed47b8cb236f00c7f5ed8/D5/S1/Depth/ContinuedFractions/KimStolarskyEventualPrefixRefutation.lean}{Lean declaration \texttt{KimStolarskyEventualPrefixRefutation.result}}
has conclusion $\neg\texttt{claim}$. This name is prefixed by
\texttt{D5.S1.Depth.ContinuedFractions}.
The proposition \texttt{claim} is the full simultaneous eventual assertion,
including the original polynomials, $m+1$ shifts, root limits and the three
original prefixes. The proof uses the squared odd term at $n=3N+2$.
The
\href{https://github.com/the-omega-institute/trureturing/blob/53c631a19a8292b44b6ed47b8cb236f00c7f5ed8/D5/S1/Depth/ContinuedFractions/QuadraticImpliesPeriodic.lean}{pinned declaration \texttt{QuadraticImpliesPeriodic.completeQuotient}}
defines the iterated inverse-fractional-part map used by \texttt{HasPrefix};
no periodicity theorem is needed for this obstruction.
The full expansions, the $L,Q$ corrections and proofs, ordering, ratio
limits and double-zero consequence have the written proofs above.
The companion \texttt{verify.py} checks the displayed polynomial identities
coefficient by coefficient, their rational sign bounds, and stated finite
ranges of continued-fraction formulas with exact integers and rational
arithmetic. These finite checks are distinct from the all-parameter
conclusions of the written proofs and from the scope of \texttt{result}.

\appendix
\section{Coefficients for the prefix bounds}\label{app:signs}
The two polynomials in~\eqref{eq:LP0}--\eqref{eq:LP1} have the following
coefficients; the row indexed by $j$ is the coefficient of $m^j$.
\begin{center}
\begin{tabular}{@{}rrr@{}}
\toprule
$j$ & $P_0$ & $P_1$ \\
\midrule
6 & 3171150 & 133974 \\
5 & 304726968 & $-335894760$ \\
4 & 9740981349 & $-50623742159$ \\
3 & 111108436490 & $-1910502470606$ \\
2 & 216415841355 & $-4118482794221$ \\
1 & 153160798296 & $-3039824512752$ \\
0 & 37000394928 & $-753289137476$ \\
\bottomrule
\end{tabular}
\end{center}
If $p_j$ are the coefficients in the last column and $m\ge M=10^6$,
then
\[
 \sum_{j=0}^{5}\frac{|p_j|}{m^{6-j}}
 \le\sum_{j=0}^{5}\frac{|p_j|}{M^{6-j}}<336,
\]
which proves the sign used above. Each identity in this appendix is a
polynomial identity obtained by clearing the explicitly positive denominators
in~\eqref{eq:LP0}--\eqref{eq:LP1}.

\begin{thebibliography}{9}
\bibitem{KS} S.-H. Kim and K. B. Stolarsky,
\emph{A balanced three or more-term generalization of Nicomachus' identity},
Integers \textbf{26} (2026), A98, pp.~1--12;
\href{https://doi.org/10.5281/zenodo.21678529}{doi:10.5281/zenodo.21678529}.
Preprint: \emph{A balanced three-term generalization of Nicomachus' identity},
\href{https://arxiv.org/abs/2511.15133v1}{arXiv:2511.15133v1} (2025),
Theorem~1, p.~5, and Remark~4, p.~6.
\bibitem{DLMF} NIST Digital Library of Mathematical Functions,
\S1.12(i), equation (1.12.1), continued-fraction notation,
\url{https://dlmf.nist.gov/1.12.i} (accessed November 17, 2026).
\end{thebibliography}
\noindent\small\copyright\ 2026 Haobo Ma and Wenlin Zhang.
Licensed under \href{https://creativecommons.org/licenses/by/4.0/}{CC BY 4.0}.
\end{document}
